(U--Th)/He dating of goethite (FeOOH) can directly date weathering, ferruginization, and
laterite formation, but interpreting goethite He dates requires knowing the conditions under
which goethite retains radiogenic helium. Vacuum step heats, the standard tool for characterizing
helium retention, convert goethite to hematite and record dehydroxylation rather than diffusion
\citep{farley2024}. Because high water pressure suppresses dehydroxylation, we use the method described
in \citep{thyer2026method} to conduct the first helium step heat
experiment on proton-irradiated goethite under water-saturated conditions. Goethite sample Win-06-01b
(Flinders Ranges, South Australia) was encapsulated in a silica glass ampoule with distilled water and heated through multiple
prograde and retrograde cycles between 170 and 250\,\textdegree C over about 40 days. Apparent
$^3$He diffusivities are approximately seven natural log units lower than those measured on the
same sample in vacuum and, beyond 7\% cumulative release, define a linear Arrhenius array with an
activation energy of \ChTwoEa{}~$\pm$~\ChTwoEaErr{}~kJ/mol and a closure temperature of
\ChTwoTc{}~$\pm$~\ChTwoTcErr{}\,\textdegree C at 10\,\textdegree C/Ma. Reanalysis of two published
vacuum step heats of the same sample shows that vacuum release follows first-order conversion to
hematite, whereas the water-saturated release shows the history dependence of diffusion and
follows multiple diffusion domain systematics. Thermodynamic and nucleation modeling calibrated on 
the vacuum data shows that the small driving force for dehydroxylation at water saturation, compared with the unbounded
driving force at the near-zero water pressure of a vacuum line, raises the 
barrier to hematite nucleation enough to suppress solid-state conversion on experimental timescales. Fluid-mediated conversion, 
observed in synthetic goethite pigment, and goethite recrystallization cannot be excluded; if either contributed, the
diffusivities are upper bounds. A four-domain model places 97\% of the $^3$He in a
single domain that is nearly quantitatively retentive at 25\,\textdegree C over billion-year
timescales. The $^4$He/$^3$He spectrum nonetheless requires natural helium loss from this domain,
which thermally activated diffusion can explain: steady cooling from 76--80\,\textdegree C at
150--175~Ma fits best, while burial of surface-formed goethite to a Late Cretaceous peak of
58--65\,\textdegree C, consistent with regional apatite fission-track data, also fits the $^4$He/$^3$He spectrum.
Water-saturated step heats enable systematic experiments
on how crystallite size, composition, and radiation damage control helium retention in goethite.
